\clearpage
\section{Protocolo de consentimiento informado}
\label{sec:consentimiento-informado}
\begingroup
\small
\setlength{\parskip}{3pt}

\noindent\textbf{Identificación.} Este protocolo corresponde al equipo MRS, integrado por Mendoza Palma Allan Jeremy, Rochina Rochina Freddy David y Sánchez Centeno Roselyn Andreina, de la asignatura Ingeniería de Requerimientos (ISR-401) de la Universidad Técnica Estatal de Quevedo. La actividad GA-TL-02 -- FORM se desarrolla bajo la docencia de Gleiston Cicerón Guerrero Ulloa, PhD.

\noindent\textbf{Propósito.} La entrevista constituye un ensayo académico de elicitación de requisitos mediante un juego de roles. La participante representa a una cajera de mostrador de un establecimiento de productos eléctricos. Sus respuestas no corresponden a un levantamiento de información realizado con trabajadores reales del sector.

\noindent\textbf{Participación y duración.} La actividad consiste en responder una entrevista virtual de aproximadamente diez minutos. Una persona del equipo conduce las preguntas y otra registra el desarrollo del ejercicio.

\noindent\textbf{Alcance de la grabación.} Se solicita autorización exclusivamente para registrar audio. Este consentimiento no autoriza la grabación de video ni de pantalla. La grabación se realizará únicamente si la participante acepta y se comprueba que el medio empleado se limita al audio autorizado.

\noindent\textbf{Voluntariedad y retiro.} La participación es voluntaria. La participante puede omitir cualquier pregunta, pedir que se detenga la grabación o retirarse de la sesión sin consecuencias. Si retira su consentimiento, el equipo dejará de utilizar el material registrado en esa sesión y gestionará su eliminación según corresponda.

\noindent\textbf{Datos personales y minimización.} La información se utilizará únicamente para la actividad académica y se tratará conforme a los principios aplicables de protección de datos personales en Ecuador~\cite{lopdp2021}. La constancia, cuando se utilice, solicitará solo nombre, fecha, modalidad del consentimiento y firma en caso de aceptación escrita. No se solicita cédula, teléfono ni correo personal.

\noindent\textbf{Confidencialidad, almacenamiento y conservación.} Los archivos de audio y las constancias que contengan datos personales no se publicarán en el repositorio de GitHub. Se almacenarán en un espacio institucional privado, con acceso restringido al equipo y al docente mediante cuentas institucionales. Su conservación se limitará a lo necesario para la evaluación académica y su eliminación se gestionará al finalizar la actividad, respetando las instrucciones institucionales y las obligaciones legales aplicables.

\noindent\textbf{Solicitud de consentimiento y constancia.} Antes de grabar, el entrevistador explica las condiciones y pregunta: ``¿Autoriza que grabemos únicamente el audio de esta entrevista académica?''. La respuesta afirmativa permite iniciar el registro de audio. Si se utiliza una constancia escrita, se completa el formulario siguiente y se conserva únicamente en el espacio institucional privado. En el repositorio público permanece el modelo sin completar.

\vspace{4pt}
\begin{center}
\renewcommand{\arraystretch}{1.85}
\begin{tabular}{@{}>{\bfseries\raggedright\arraybackslash}p{0.30\linewidth}>{\raggedright\arraybackslash}p{0.61\linewidth}@{}}
Nombre de la participante & \rule[-2pt]{0.95\linewidth}{0.4pt} \\
Fecha & \rule[-2pt]{0.95\linewidth}{0.4pt} \\
Modalidad de consentimiento & Verbal\quad \fbox{\rule{0pt}{1.6ex}\rule{1.6ex}{0pt}}\quad Escrita\quad \fbox{\rule{0pt}{1.6ex}\rule{1.6ex}{0pt}} \\
Firma (solo si es escrito) & \rule[-2pt]{0.95\linewidth}{0.4pt} \\
\end{tabular}
\end{center}
\endgroup
